\subsection{Workflow Interruptions and Safe Next Actions}
\label{sec:adoption-workflow}

The Agent Adapter and product UI expose the existing result and lifecycle
semantics faithfully. Table~\ref{tab:adoption-handling} is integration guidance,
not a new wire error or state profile. Each protected task shows the exact
action, a policy reason safe to disclose, its proposal/attempt reference,
current status, responsible reviewer/operator and next supported action.
Deduplicate review prompts by stable proposal/operation identity. Editing a
reviewed action creates a new proposal and admission; a terminal confirmation
cannot substitute for a required hardware-review obligation.

{\small
\begin{longtable}{@{}L{0.22\linewidth}L{0.42\linewidth}L{0.31\linewidth}@{}}
\caption{Existing results mapped to workflow handling and user communication.}
\label{tab:adoption-handling}\\
\toprule
\textbf{Result or Lifecycle State} & \textbf{Workflow Response} & \textbf{User Status / Accountable Owner}\\
\midrule
\endfirsthead
\toprule
\textbf{Result or Lifecycle State} & \textbf{Workflow Response} & \textbf{User Status / Accountable Owner}\\
\midrule
\endhead
\bottomrule
\endfoot
\nolinkurl{NEEDS_EVIDENCE} & Pause the unchanged proposal; validate the trusted remediation path and obtain only the required evidence. Resume exact admission rather than replanning silently. & Awaiting named evidence/review, with expiry and a status reference. Application/product owner and review provider.\\
\addlinespace
\nolinkurl{DENIED}; identity/grant mismatch & Stop protected dispatch. Route policy changes through authorized administration; investigate identity/binding errors. Do not repeatedly ask another reviewer for the same refused action. & Denied or binding error, safe reason and escalation path. Policy/security owner.\\
\addlinespace
Stale policy/evidence; \nolinkurl{STATE_GUARD_VIOLATION} & Refresh and revalidate through admission. Inspect existing root/attempt state before issuing replacement authority; changed bound context or parameters require the appropriate new proposal. & Action needs revalidation, not completed. Admission/resource owner.\\
\addlinespace
\nolinkurl{ROOT_UNAVAILABLE}; \nolinkurl{NONCE_CONSUMED} & Query the existing authority/attempt. A new RPC ID, nonce or automatic tool retry cannot create an additional allowance. & Existing operation or authority issue; status lookup and operator escalation. Platform owner.\\
\addlinespace
\texttt{dispatched} / \texttt{pending} & Query the correlated target operation under its completion contract; do not redispatch to obtain a faster response. & Accepted or executing, with no premature completion claim. Resource adapter/SRE.\\
\addlinespace
\texttt{unknown}; \nolinkurl{OUTCOME_UNKNOWN} & Retain consumption and any reservation; reconcile authoritative evidence and escalate at the declared deadline. Another action does not clear uncertainty. & Outcome unknown, last confirmed fact and incident reference. Resource owner/SRE.\\
\addlinespace
\texttt{committed}; \texttt{failed} / \texttt{compensation\_failed} & Report the receipt's actual effect scope. Failure does not imply zero effect or refund; any supported compensation needs fresh admission and known preconditions. & Confirmed effect or failure with possible partial effect. Resource/product owner.\\
\addlinespace
\texttt{not\_dispatched}; \texttt{fenced\_no\_effect} & Preserve the verified refusal/resolution; \texttt{not\_dispatched} alone does not restore an allowance. Authoritative no-effect evidence plus the required late-effect barrier can support restoration under issuer conservation rules. & No dispatch or fenced no effect, with supporting evidence. Gateway/resource owner.\\
\end{longtable}
}

Before claim, UI cancellation may withdraw a proposal or request atomic
retirement of unclaimed authority under the authority contract. For claimed
or dispatched work, cancellation stops future work only under the declared
target contract: it does not prove no effect, release a reservation or undo a
commit. Without authoritative resolution and any required late-effect barrier,
the attempt remains pending or unknown. The product owner explains this
distinction before launch.

\subsection{Operate, Upgrade, and Exit Without Reopening Authority}
\label{sec:adoption-operations}

During an outage, permitted planning and workspace work can continue within
their policy. New protected dispatch stops when required admission, claim or
credential authority is unavailable. An Outcome Service outage may buffer
already supported records in durable outboxes under the declared contract;
if a required transition record cannot be durably retained, that transition
must stop. Disconnected operation requires a supported local authority/target
profile preserving custody, freshness, atomic claim and recovery. Cached policy
or admission alone is insufficient.

The deployment owner rehearses cutover, abort, upgrade, restore and supplier
exit as one authority-preserving transition:
\begin{enumerate}
\item Quiesce new issuance and dispatch; establish the resource's fencing or
late-effect barrier where required. Inventory every granted, claimed, pending
and unknown attempt. Stopping the controller does not cancel target work.
\item Reconcile or quarantine in-flight operations. Retain consumed roots,
generations, reservations, replay state, dispatch intents and outboxes; a stale
restore or a fresh empty ledger cannot make old authority reusable.
\item Preserve original signed records, trusted key/version histories and
descriptor/policy bindings. Validate the new runtime/profile and key transition;
incompatible signed records are not silently rewritten or re-signed as proof
of the old history.
\item Run the same scope-specific refusal, bypass, concurrency, lost-response
and restore cases before resuming. Separately authorized emergency operators
remain declared writers. Returning equivalent ambient credentials to an agent
changes or ends the protected scope and requires explicit owner acceptance.
\end{enumerate}

The acceptance record names dependency patching/provenance, key rotation,
outage backlog, recovery deadlines and on-call ownership. A failed pilot can
disable the protected feature while retaining permitted work and a separately
governed operator procedure. This makes postponement practical without treating
an unsafe fallback as a successful migration.

\subsection{Buyer Control and Supplier Evidence Record}
\label{sec:buyer-control-record}

All buyers need a deployment-specific evidence record. Government and
sovereign procurement teams additionally use it to state who controls the
deployment and its dependencies; local hosting alone does not establish that
control. Name the operators and administrators of policy, grant signing,
workload enrollment, Evidence Providers, credential brokers, optional grant
stores/resolvers, gateways,
authority state and retained outcomes. Record which external model, judge,
identity, signing or support service can observe, alter or disable the chosen
operation, and whether the buyer can operate the required path independently.
Declare provider independence and conflicts of interest rather than inferring
them from multiple signatures.

Inventory the sensitive canonical proposals, identity, evidence, review and
receipt records: access, storage/transfer location, retention, permitted
disclosure and export. Redacted views retain links to the protected originals;
altering a signed payload breaks its signature and digest bindings. Specify
who may verify the retained chain, what records verification requires, and how
originals/key histories survive supplier exit.

\small
\begin{table}[!ht]
\begin{tabularx}{\linewidth}{@{}L{0.22\linewidth}L{0.34\linewidth}X@{}}
\toprule
\textbf{Acceptance Topic} & \textbf{Evidence Supplied Here} & \textbf{Evidence Required from Deployment}\\
\midrule
Contract exchange & Closed schemas, signed vectors and independent canonicalization/signature checks. & Frozen versions and independently implemented peers exchanging the selected profile.\\
Protected boundary and review & Proxy identity/claim tests; mock guards; software-key WebAuthn mechanics. & Actual enrollment, hardware assurance when required, native bypass/custody and selected-target guard/outcome tests.\\
Recovery and operation & Bounded claim races and mock lost-response reconciliation; architecture acceptance cases. & Real fault/restore, key transition, emergency operation, data controls, support and safe supplier-exit rehearsal.\\
Adoption value & Proposed work packages, friction responses and measurement method. & Workload-specific user/operating cost, accepted guarantee limits and signed-off expansion decision.\\
\bottomrule
\end{tabularx}
\caption{Procurement evidence split. Section~\ref{sec:executed-checks} states the
scope of supplied checks; Section~\ref{sec:acceptance} supplies deployment tests.
This is an acceptance input, not a certification or jurisdiction-specific
accreditation claim.}
\label{tab:buyer-evidence}
\end{table}
\normalsize

The resource/public-service owner accepts policy correctness and the residual
limits of B/C adapters, external writers and uncertain effects. The buyer
selects any legal, national-assurance or accreditation requirements separately
and obtains evidence for its chosen profile. No sovereign-operation,
production-resilience or supplier-portability result is established by this
whitepaper.
